% The wall dictionary.  Split out of the fixed-count transport section so that
% the sections that consume it --- the spike, the vertex-on-edge laws --- come
% after it in document order and not before.
\section{Relative general position and the wall dictionary}\label{sec:dictionary}

A path between generic polygons meets the discriminant, and every law proved
later is a statement about one kind of crossing. This section fixes the regular
locus, replaces an arbitrary path by a finite factor-simple path, and proves
that each relevant hit carries exactly one of five bundles. Nothing here
consumes a wall law; the sections that prove those laws consume this one.

\subsection{The regular locus}


\begin{lemma}[the regular locus is open, and generic polygons lie in it]
\label{lem:regularopen}
$\mathcal R_n$ is open in $(\mathbb R^2)^n$, and every generic polygon
(Definition~\ref{def:generic}(A)) lies in $\mathcal R_n$. Consequently a compact
path in $\mathcal R_n$ has a positive minimum edge length and positive distance
from the complement.
\end{lemma}
\begin{proof}
Write $B_i=\{d_i=0\}$ and
$C_i=\{d_{i+1}=-\lambda d_i\text{ for some }\lambda>0\}$. Each $B_i$ is closed.
The union $B_i\cup B_{i+1}\cup C_i$ is closed: let $P^{(k)}$ be a sequence in it
converging to $P$, and suppose $P\notin B_i\cup B_{i+1}$, so $d_i,d_{i+1}\neq0$
at $P$. Then for large $k$ the same holds at $P^{(k)}$, so $P^{(k)}\in C_i$ and
$d_{i+1}^{(k)}=-\lambda_kd_i^{(k)}$ with $\lambda_k>0$; taking norms,
$\lambda_k=|d^{(k)}_{i+1}|/|d^{(k)}_i|\to|d_{i+1}|/|d_i|>0$, and passing to the
limit gives $d_{i+1}=-\lambda d_i$ with $\lambda>0$, i.e.\ $P\in C_i$. The set
$B_i\cup C_i$ alone need \emph{not} be closed --- with $d_i=(1,0)$ fixed and
$d_{i+1}=(-1/k,0)$, every term lies in $C_i$ while the limit lies only in
$B_{i+1}$ --- which is why the triple union is what the paragraph proves
closed. Since the indices are cyclic,
$\bigcup_i(B_i\cup C_i)=\bigcup_i(B_i\cup B_{i+1}\cup C_i)$, a finite union of
the closed sets just proved closed; that union is the complement of
$\mathcal R_n$, so $\mathcal R_n$ is open. An earlier revision called
$B_i\cup C_i$ closed on the way. A generic polygon has
$\det(d_{i-1},d_i)\neq0$ for all $i$ by (G1), so no consecutive pair is
proportional at all. The last sentence is compactness.
\end{proof}




\subsection{Finite relative general position, and which bundles occur}

\begin{lemma}[canonical factors: definition and finite relative general position]\label{lem:relgp}
Fix $n\geq3$.
On $\mathbb R^{2n}$ define, for three distinct vertex indices $a,b,c$, for two
\emph{remote} edges $e\neq f$, and for three pairwise remote edges $e,f,g$,
\[
   \Delta_{abc}=\det\bigl(p_b-p_a,\ p_c-p_a\bigr),\qquad
   H_{e,f}=\det(d_e,d_f),\qquad
   T_{e,f,g}=\det\bigl(\ell_e,\ell_f,\ell_g\bigr),
\]
where $\ell$ is the line-coefficient row of Definition~\ref{def:guarded}. A
\emph{canonical factor} is a polynomial in one of these three families with
indices as restricted here. Two canonical factors are \emph{associates} if one
is a nonzero real multiple of the other; the associate classes are what the
counting below uses, and an earlier revision spoke of canonical factors without
ever defining either the list or the equivalence.

The index restrictions are not decoration. With $a$, $b$, $c$ not distinct,
$\Delta$ is the zero polynomial; likewise $\det(d_e,d_e)=0$, and
$\det(d_e,p_i-p_e)=0$ when $p_i$ is an endpoint of $e$. An earlier revision
printed an unrestricted list and called every member of it irreducible, which
is false at exactly those indices.

The turn and the vertex-on-line determinant are one family and not two:
\[
   \det\bigl(p_b-p_a,\ p_c-p_b\bigr)=\det\bigl(p_b-p_a,\ p_c-p_a\bigr)
   =\Delta_{abc},
\]
identically, so the turn at $b$ of the two edges $ab$, $bc$ \emph{is} the
determinant that puts $c$ on the line of $ab$. Every (G1) and every (G2) is a
$\Delta$, and so is every direction determinant of two \emph{consecutive}
edges; $H$ is reserved for remote pairs, where four distinct vertices occur.

\begin{enumerate}
\item[(A1)] Every member of $\mathcal G$ is, up to sign, a canonical factor
as defined above or a product of two: $\mathrm{G1}_j=\Delta_{(j-1)j(j+1)}$,
$\mathrm{G2}_{e,i}=\Delta_{e_0e_1i}$ with $e=(e_0,e_1)$ and $i$ not an endpoint
of $e$, $\mathrm{G5}_{e,f}=H_{e,f}$, and $\mathrm{G3}_{e,f,g}=T_{e,f,g}$; the
order predicate $\mathrm{G4}_{e;f,g}$ equals $-T_{e,f,g}$ when $f$ and $g$ are
remote to each other, and $-\Delta_{e_0e_1M}\Delta_{f_0Mg_1}$ when they meet at
the vertex $M$, which is a product of two canonical factors and the only
member that is not a single one.
\item[(A2)] Every canonical factor is irreducible in $\mathbb R[x_1,\dots,x_{2n}]$.
\item[(A3)] Two canonical factors are associates if and only if they have the same
variable support and one is a signed reindexing of the other; explicitly, the
associate class of $\Delta_{abc}$ is the six orderings of $\{a,b,c\}$ with the
sign of the permutation, that of $H_{e,f}$ is $\{H_{e,f},-H_{e,f}\}$, and that
of $T_{e,f,g}$ is its six row orderings with sign.
\item[(B)] Let $\mathcal C$ contain one representative of each associate class of the
canonical factors occurring in members of $\mathcal G$. Let
$\gamma\colon[0,1]\to\mathcal R_n$ be continuous with generic endpoints and
let $\delta>0$. There is a continuous piecewise-affine path
$\tilde\gamma\colon[0,1]\to\mathcal R_n$, with the same endpoints and within
$\delta$ of $\gamma$, whose non-generic parameters form a finite set. Every
such parameter $t_0$ lies in the interior of a nonconstant segment on which one
scalar vertex coordinate $x_j$ moves, and has this certificate:
\begin{enumerate}
\item[(B1)] exactly one $F\in\mathcal C$ vanishes at
$x=\tilde\gamma(t_0)$;
\item[(B2)] $(F\circ\tilde\gamma)'(t_0)\neq0$;
\item[(B3)] no two vertices coincide at $x$.
\end{enumerate}
Every junction outside the two endpoint collars has all canonical factors
nonzero and has both scalar-coordinate differences nonzero for every pair of
vertices. The endpoint collars lie in $\mathcal U_n$. Thus corners of the path
are wall-free, and every wall germ is an affine one-parameter family with one
simple canonical-factor root.
\end{enumerate}
\end{lemma}
\begin{proof}
For (A1), use Definition~\ref{def:guarded} together with the displays above and the
factorization of a concurrency determinant whose last two edges are
consecutive.

For (A2), consider $\Delta$ and $H$: the quadratic form $q(u,w)=u_1w_2-u_2w_1$ has rank
four, and a reducible quadratic form is a product of two linear forms and so has
rank at most two; hence $q$ is irreducible. Both are $q$ pulled back along a
surjective linear map of the vertex coordinates --- $(p_b-p_a,\,p_c-p_a)$ in one
case, $(d_e,d_f)$ in the other, and both maps are surjective precisely because
the index restrictions make the two arguments independent. A pullback $q\circ L$
along a surjective linear $L$ is again irreducible: a factorization $q\circ L=AB$
has both factors independent of $\ker L$, because a divisor of a polynomial of
degree zero in a variable has degree zero in it, so $A=A'\circ L$ and
$B=B'\circ L$ and $q=A'B'$.

For $T$: it has degree one in each of the six two-vector blocks
$p_e,d_e,p_f,d_f,p_g,d_g$, which are independent coordinates because the edges
are pairwise remote. Suppose $T=AB$ with both factors nonconstant. Since $T$ has
degree one in $d_g$, one factor, say $B$, is free of $d_g$. Expanding the
determinant along the row of $g$ and collecting,
\[
   T=d_{g,y}\bigl(p_{g,x}M_C-M_A\bigr)+d_{g,x}\bigl(M_B-p_{g,y}M_C\bigr),
\]
where $M_A,M_B,M_C$ are the $2\times2$ cofactors of that row and involve only
the blocks of $e$ and $f$. As $B$ is free of $d_g$ it divides both bracketed
polynomials; the first is free of $p_{g,y}$ and the second is free of
$p_{g,x}$, so $B$ is free of both, hence of every block of $g$. Comparing the
coefficients of the four distinct monomials $d_{g,y}p_{g,x}$, $d_{g,y}$,
$d_{g,x}$, $d_{g,x}p_{g,y}$ then gives $B\mid M_A$, $B\mid M_B$ and
$B\mid M_C$. But $M_C=\det(d_e,d_f)=H_{e,f}$, which is irreducible by the
previous paragraph, so $B$ is a unit or an associate of $H_{e,f}$; and
$H_{e,f}$ does not divide $T$, since setting $d_e=d_f=(1,0)$ makes $H_{e,f}$
vanish while $T$ becomes $-d_{g,y}(p_{e,y}-p_{f,y})$, not the zero polynomial.
So $B$ is a unit, contrary to assumption.
For (A3), first observe that associates have the same zero set and the same variable support, so it is
enough to consider two canonical factors of the same family on the same
support. For $\Delta_{abc}$ the six orderings of $a,b,c$ give the same polynomial up to
the sign $\sgn\sigma$ of the permutation $\sigma$ relating them, and nothing
else: $\Delta_{\sigma(a)\sigma(b)\sigma(c)}=\sgn(\sigma)\,\Delta_{abc}$,
a determinant being alternating in its rows. For $T_{e,f,g}$ the same holds of
the six orderings of the three line-coefficient rows,
$T_{\sigma(e)\sigma(f)\sigma(g)}=\sgn(\sigma)\,T_{e,f,g}$. So two orderings of
one triple are associates by a unit $\pm1$, which is what this clause needs; an
earlier revision said "the sign of the permutation" without binding $\sigma$ or
displaying either identity. Two canonical
factors on \emph{different} supports are never associates: a vertex coordinate
occurring in one and not in the other cannot be matched by a scalar.

Two same-support pairs remain, and they are the only ones. The support of a
remote edge pair $\{e,f\}$ is $\{e,e+1\}\sqcup\{f,f+1\}$, two disjoint blocks of
two cyclically adjacent indices. If a second remote pair has the same support,
that four-element set splits into two adjacent blocks in a second way, so some
index is blocked with its \emph{other} neighbour: $f=e+2$ and, applying the same
to the other block, $e=f+2$. Then $e=e+4$ in $\mathbb Z_n$, so $n$ divides $4$,
and $n\geq3$ leaves $n=4$. The same count settles triples: the support of a
pairwise remote triple is three disjoint adjacent blocks, a second such
decomposition pairs every index with its other neighbour, and following the
chain $a\to a+1\to a+2\to\cdots$ around returns to $a$ after six steps, so
$n=6$.

At $n=4$ the two pairs are
$H_{0,2}=\det(p_1-p_0,p_3-p_2)$ and $H_{1,3}=\det(p_2-p_1,p_0-p_3)$: the
monomial $x_0y_2$ occurs in the first with coefficient $1$ and in the second
with coefficient $-1$, while $x_0y_1$ occurs in the second and not in the first,
so no scalar carries one to the other. At $n=6$ the two triples are the
alternating ones, and the hexagon
$p_0=(1,0)$, $p_1=(2,0)$, $p_2=(0,1)$, $p_3=(0,2)$, $p_4=(1,1)$, $p_5=(2,2)$
gives $T_{0,2,4}=0$ and $T_{1,3,5}=2$, so those two are not associates either.


For (B), put $N=2n$ and write the scalar coordinates as $x_1,\ldots,x_N$. The set
$\mathcal C$ is finite by (A1). We first print the
one-variable algebra used to choose the path.

Fix $F\in\mathcal C$ depending on $x_j$, and put
$K_j=\operatorname{Frac}(\mathbb R[x_1,\ldots,\widehat{x_j},\ldots,x_N])$.
Irreducibility of $F$ in the full polynomial ring
(A2) and Gauss's lemma make it irreducible in
$K_j[x_j]$. Characteristic zero and dependence on $x_j$ give
$\gcd(F,\partial_jF)=1$ there. A Bezout identity, after clearing its finitely
many denominators, therefore has the form
\[
 A_{F,j}F+B_{F,j}\partial_jF=D_{F,j},
 \qquad 0\neq D_{F,j}\in
 \mathbb R[x_1,\ldots,\widehat{x_j},\ldots,x_N].
\]
If nonassociate $F,G\in\mathcal C$ both depend on $x_j$, they are coprime in
$K_j[x_j]$: otherwise their primitive irreducible representatives would become
associates over $K_j$, and clearing the scalar would make them associates in
the full ring. Clearing a second Bezout identity gives
\[
 A_{FG,j}F+B_{FG,j}G=R_{FG,j},
 \qquad 0\neq R_{FG,j}\in
 \mathbb R[x_1,\ldots,\widehat{x_j},\ldots,x_N].
\]
At a specialization where $D_{F,j}\neq0$, the univariate specialization of
$F$ has no repeated root; where $R_{FG,j}\neq0$, the specializations of $F$
and $G$ have no common root. A factor independent of $x_j$ is constant on that
coordinate segment and is kept nonzero at its initial junction. This includes
the two exceptional same-support comparisons: the $n=4$ pair of $H$ factors
and the $n=6$ alternating pair of $T$ factors are nonassociate by the explicit
monomial and hexagon tests in (A3), so their Bezout
certificates occur in the same finite list.

We now choose the whole path at once, including its parametrization. Use the
same norm as in the asserted $\delta$-bound. Compactness of $\gamma([0,1])$
and openness of $\mathcal R_n$ (Lemma~\ref{lem:regularopen}) give an open
cover of $\gamma([0,1])$ by axis-aligned cubes whose closures lie in
$\mathcal R_n$ and whose diameters are less than $\delta$. At the two
endpoints choose the cover cubes inside $\mathcal U_n$, using openness of
$\mathcal U_n$ (Proposition~\ref{prop:chamberinv}(i)). Uniform continuity of $\gamma$
(equivalently, a Lebesgue-number subdivision of the pulled-back cover) gives
a partition
\[
 0=t_0<t_1<\cdots<t_m=1
\]
and cubes $Q_0,\ldots,Q_{m-1}$ from this cover such that
$\gamma([t_i,t_{i+1}])\subset Q_i$; refine at the ends so that
$Q_0,Q_{m-1}\subset\mathcal U_n$. In particular
$\gamma(t_i)\in Q_{i-1}\cap Q_i$, so every consecutive overlap is nonempty.
Put $q_0=\gamma(0)$ and $q_m=\gamma(1)$, and choose each internal waypoint
$q_i$ in $Q_{i-1}\cap Q_i$. Then $q_i,q_{i+1}\in Q_i$, so every
coordinate-order hybrid between them remains in $Q_i$.

Call the intervals with $1\leq i\leq m-2$ central. Their waypoints are chosen
jointly in the nonempty open product of the overlaps. Impose the following
finite list of nonvanishings for every central
coordinate segment and every one of its hybrids: $F$ at the hybrid for each
$F\in\mathcal C$; the corresponding $D_{F,j}$ and $R_{FG,j}$ on the fixed
coordinates of the segment; the coordinate step
$q_{i+1,j}-q_{i,j}$; and, at every hybrid, both coordinate differences for
every pair of vertices. Each condition pulls back to a nonzero polynomial on
the waypoint product. Indeed, a hybrid or fixed-coordinate map selects distinct
scalar coordinates from the two adjacent variable waypoints, so its image
contains an open coordinate box; a nonzero polynomial cannot vanish on that
box. The step and vertex-difference conditions are visibly nonzero linear
polynomials. Their finite product is nonzero because the polynomial ring is an
integral domain, and a nonzero polynomial cannot vanish on a nonempty open set
(all Taylor coefficients at an interior point would then vanish). Hence one
joint waypoint choice satisfies every condition. Sequential choice is not used.

On each original interval $[t_i,t_{i+1}]$, join $q_i$ to $q_{i+1}$ by
changing the $N$ scalar coordinates in the fixed order and parametrizing the
resulting coordinate legs inside that same interval (constant legs are simply
omitted). The first and last such pieces lie in
$Q_0,Q_{m-1}\subset\mathcal U_n$ and therefore contain no wall, even if an
inactive canonical factor vanishes there. For every $t\in[t_i,t_{i+1}]$, both
$\tilde\gamma(t)$ and $\gamma(t)$ lie in $Q_i$. Consequently
\[
 \|\tilde\gamma(t)-\gamma(t)\|<\operatorname{diam}(Q_i)<\delta,
\]
which is the required pointwise, hence uniform, bound. The central path lies
in $\mathcal R_n$; its junctions are factor-free, and every leg is nonconstant.
Along a leg, the Bezout certificates make each canonical-factor root simple
and prevent roots of two nonassociate classes from coinciding. Each
restriction is a nonzero univariate polynomial, so the union of their root
sets over finitely many legs is finite.

Finally, vertices cannot collide. On a leg moving an $x$-coordinate, every
relevant $y$-difference is the fixed nonzero value imposed at its two hybrid
ends; on a leg moving a $y$-coordinate, the corresponding $x$-difference is
fixed and nonzero; pairs not involving the moving vertex do not change. A
non-generic parameter is a zero of a relevant member, hence of one of its one
or two canonical factors (A1); the root audit gives
the unique $F$, its simple derivative, and the collision exclusion claimed.
\end{proof}

\begin{remark}[bundles, not single predicates]\label{rem:bundles}
Two mistakes in an earlier general-position route are worth keeping visible.
It treated any two simultaneously vanishing members as independent conditions,
which is false for every bundle below; and it asserted that every active (G4)
tie is a three-remote-edge concurrence, missing the incident-edge tie of the
bigon. The exact witness is the remote edge $(0,0)\to(4,0)$ with
$M=(2,t)$, $p=(1,1)$ and $q=(3,1)$: there
$\mathrm{G2}_{e,M}=4t$ and $\mathrm{G4}_{e;f,g}=-8t(1-t)$.
\end{remark}

\begin{remark}[where bundle transversality is consumed]\label{rem:bundlescope}
At a Reidemeister-III hit the induction verifies event transversality directly
from the common simple factor supplied by Lemma~\ref{lem:dictionary}; at a
silent hit no sign change is consumed. Thus the fixed-count transport step in the induction
consumes the dictionary's simple-zero conclusion only for flat, bigon and
sliding bundles; the induction treats Reidemeister III separately.
\end{remark}

\begin{lemma}[inactive-zero exclusions and the exhaustive wall dictionary]\label{lem:dictionary}
\begin{enumerate}
\item[(A)] Let $t\mapsto P(t)$ be an event and let $Z$ be its zero set.
\begin{enumerate}
\item[(a)] If $\mathrm{G3}_{e,f,g}$ vanishes at $t=0$, the three lines are not
all equal, and their common point fails to lie on at least one of the three
segments, then $\mathrm{G3}_{e,f,g}\notin Z$. (When the three lines coincide
there is no common point to speak of and the hypothesis says nothing; that case
is excluded here and is not used below.)
\item[(b)] If $\mathrm{G5}_{e,f}$ vanishes at $t=0$ and the two closed segments
$e(0)$ and $f(0)$ are disjoint, then $\mathrm{G5}_{e,f}\notin Z$.
\end{enumerate}
\item[(B)] Let $\tilde\gamma$ be a path supplied by Lemma~\ref{lem:relgp}(B), let $t_0$ be
one of its finitely many non-generic parameters, and put
$x=\tilde\gamma(t_0)$. Define
\[
 \mathcal B(x)=\{g\in\mathcal G:g(x)=0\text{ and }g\text{ is relevant at }
 \tilde\gamma(t)\text{ for some }t\neq t_0\text{ in every neighbourhood of }t_0\}.
\]
On an interval containing no other non-generic parameter this is exactly the
zero set of the event of Definition~\ref{def:event}. There is one canonical
factor $F$ with $F(x)=0$ and $\mathrm dF(x)\neq0$ such that every
$g\in\mathcal B(x)$ has $g=F B_g$ with $B_g(x)\neq0$. In particular all
members share one smooth branch and their restrictions to the path have simple
zeros, so every member changes sign at $t_0$. The bundle is exactly one of:
\begin{enumerate}
\item[(i)] \emph{Flat}: $\{\mathrm{G1}_j,
\mathrm{G2}_{(j-1,j),j+1},\mathrm{G2}_{(j,j+1),j-1}\}$ with $p_j$ strictly
between its neighbours; wall (S3).
\item[(ii)] \emph{Vertex-on-edge, bigon}:
$\{\mathrm{G2}_{(a,b),M},\mathrm{G4}\langle(a,b);e_{M-1},e_M\rangle\}$,
the foot in the segment interior and the two neighbours of $M$ strictly on one
side of the remote line; wall (S7).
\item[(iii)] \emph{Vertex-on-edge, sliding}: the singleton
$\{\mathrm{G2}_{(a,b),M}\}$, the foot in the segment interior and the two
neighbours strictly on opposite sides; wall (S7).
\item[(iv)] \emph{Reidemeister III}: $\{\mathrm{G3}_{e,f,g}\}$ together with
the three (G4) members whose ties it forces, for three pairwise remote edges
concurrent in the interior of all three; wall (S2).
\item[(v)] \emph{Silent}: a (G2) vanishing with its point off the segment,
together with whatever members it forces, silent in the sense of
Definition~\ref{def:silent}.
\end{enumerate}
The exterior consecutive-collinearity branch is the cusp and lies outside
$\mathcal R_n$. Remote-edge overlap, vertex coincidence, four-edge
concurrence and two independent bundles cannot occur at a certified hit.
\end{enumerate}
\end{lemma}
\begin{proof}
\emph{Clause (A).} \emph{Clause (a).} Say the common point $q$ misses
the segment $g$. Activation asks that all
three pairs cross (Definition~\ref{def:guarded}), so it is enough to exhibit
\emph{one} pair that does not cross at $t=0$ and does not cross for $|t|$ small,
and the pair must be chosen with care: an earlier revision took $(e,g)$ and said
"the lines of $e$ and $g$ meet only at $q$", which is false when those two lines
coincide.

Choose the pair as follows. The three lines are not all equal, so at least one
of $e,f$ has a line different from $g$'s; call it $h$, so that the lines of $h$
and $g$ are distinct and meet exactly at $q$. Since $q$ misses the segment $g$,
the edges $h$ and $g$ do not cross at $t=0$. For nearby $t$: if the lines of $h$
and $g$ are not parallel their intersection point is a continuous function of
the vertices, and at $t=0$ it lies at positive distance from the compact segment
$g$, so it misses $g$ for all $|t|$ small; if they are parallel --- which at
$t=0$ they are not, distinct lines through $q$ --- they can become so only by
first ceasing to meet near $g$, and in any case two parallel edges meet at most
in a collinear overlap and do not cross transversally. Either way $h$ and $g$ do
not cross for $|t|$ small, so $\mathrm{G3}_{e,f,g}$ is inactive there and
relevant at no $t\neq0$ near $0$. By Definition~\ref{def:event} it is not in
$Z$.

\emph{Clause (b).} Disjoint compact sets stay disjoint under a small
perturbation: the
distance $\min\{|u-v|:u\in e(t),v\in f(t)\}$ is a continuous function of $t$,
positive at $t=0$, hence positive for $|t|$ small. So $e$ and $f$ do not meet,
in particular do not cross, for those $t$; $\mathrm{G5}_{e,f}$ is inactive
there and relevant at no such $t\neq0$, so it is not in $Z$.

\emph{Clause (B).} Choose an interval $I$ about $t_0$ containing no other non-generic parameter.
Every point of $I\setminus\{t_0\}$ is generic, so the restricted path is an
event. Activation is constant on each punctured side because it is decided by
strict sign conditions in unconditional (G2) members. This proves the asserted
equality between $\mathcal B(x)$ and the event zero set.

\emph{Event-cell check.} Fix a vertex $M$ and a remote edge $e$. Put
$s_{e,M}=\langle p_M-p_e,d_e\rangle$. Its three strict incidence regimes are
$s_{e,M}<0$, $0<s_{e,M}<\lVert d_e\rVert^2$ and
$s_{e,M}>\lVert d_e\rVert^2$, according as the perpendicular foot lies before,
inside or after the segment; the outer two are the silent regime. The two
strict side regimes are distinguished by the sign of
$\mathrm{G2}_{e,M-1}\mathrm{G2}_{e,M+1}$: positive puts the neighbours on one
side of the remote line and negative puts them on opposite sides. Finally, for
each $g\in\mathcal G$, its activation locus is the finite conjunction of
strict guarded inequalities of Definition~\ref{def:guarded} (all of
$\mathbb R^{2n}$ when $g$ is unconditional), and its two sign regimes are
$g>0$ and $g<0$. All these regimes are semialgebraic, being cut out by finitely
many polynomial equalities and strict inequalities. This is their only use.

By Lemma~\ref{lem:relgp}(A1), every member of $\mathcal G$ is, up to
sign, one canonical factor or a product of two. Lemma~\ref{lem:relgp}(B1) says that
exactly one associate class vanishes at $x$. Thus a vanishing single-factor
member is a unit multiple of its representative $F$; for a product exactly one
factor is associate to $F$ and the other is nonzero at $x$. Hence every member
of $\mathcal B(x)$ is $F B_g$ with $B_g(x)\neq0$. The simple derivative of
$F\circ\tilde\gamma$ implies $\mathrm dF(x)\neq0$ and makes every restricted
member change sign with the same branch.

It remains to classify the possible defining family. A (G1) zero inside
$\mathcal R_n$ is positive collinearity, so the middle vertex lies between its
neighbours; its two (G2) companions are identically the same determinant. This
is (i). Negative collinearity is outside $\mathcal R_n$ and is the cusp.

For a (G2) zero, first remove the adjacent case: if $M=a-1$ or $M=b+1$, the
member is the corresponding (G1), hence (i). Otherwise the perpendicular foot
cannot be an endpoint of the remote segment, because at a zero that would make
$p_M$ equal to a remote-edge endpoint, a vertex collision excluded by
Lemma~\ref{lem:relgp}(B3). Nor can a neighbour of $M$ lie on the remote line: its
vertex-on-line factor is nonassociate to the defining factor outside the
already separated adjacent case, so two classes would vanish. The incidence
and side regimes in the event-cell check above are therefore strict. With
the foot inside the segment, neighbours on one side give the bigon bundle (ii):
the order member is active on the two-crossing side and its consecutive-edge
factorization has the incidence as its unique vanishing factor. Neighbours on
opposite sides give the sliding singleton (iii), since the order member is
active on neither side. With the foot off the segment, clause~(A)
leaves the silent case (v).

For three pairwise remote edges, the concurrency determinant and the three
order predicates satisfy, for a fixed ordering,
\[
 \mathrm{G4}_{e;f,g}=-\mathrm{G3}_{e,f,g},\qquad
 \mathrm{G4}_{f;e,g}=+\mathrm{G3}_{e,f,g},\qquad
 \mathrm{G4}_{g;e,f}=-\mathrm{G3}_{e,f,g}.
\]
When the three crossings are active their common point is interior to all three
segments, and this is (iv). If the last two edges of a (G4) are consecutive at
$M$, its polynomial instead factors as the vertex-on-line factor
$\Delta_{e_0e_1M}$ times the turn at $M$; the unique-factor certificate assigns
the hit to the corresponding (G2) or (G1) case. A punctured-side-relevant (G5) contributes no new case. Suppose it
vanishes at $x$. If some (G1) or (G2) member also vanishes there, the preceding
flat or vertex-on-edge case owns the hit, and the unique-factor certificate
leaves no independent (G5) family. Otherwise all (G1) and (G2) members are
nonzero at $x$. Their strict signs persist to $x$; an endpoint limit would force
a (G2) zero, so the pair is active at $x$. Proposition~\ref{prop:fidelity}(i)
then contradicts the (G5) vanishing.

Finally, each excluded coincidence violates a printed part of the certificate.
Remote-edge overlap forces two nonassociate vertex-on-line factors. Four-edge
concurrence forces two nonassociate concurrency factors; the possible
same-support $n=6$ pair is nonassociate by Lemma~\ref{lem:relgp}(A3).
Two independent bundles likewise force two nonassociate classes. A vertex
coincidence is excluded separately by Lemma~\ref{lem:relgp}(B3). These
exhaust the guarded families and the event-cell boundaries, proving the list.
\end{proof}

